\section{问题三：连续通信约束下的运输与中继协同}
\subsection{双向链路预算与可用性}
固定网关G01位于O01，通信端点海拔由地面海拔与天线离地高度确定；运输端点随爬升、巡航、下降和交接阶段变化。接收门限为灵敏度与衰落裕量之和，附件中相应值为
\begin{equation}P_{th}=P_{sens}+M=-98+8=-90\ \mathrm{dBm}.\end{equation}
对通信方向$a\to b$，允许传播损耗为
\begin{equation}
L_{a\to b}^{max}=P_t^a+G_t^a+G_r^b-L_{sys}-P_{th}^b,\qquad
L_{ab}^{max}=\min(L_{a\to b}^{max},L_{b\to a}^{max}).
\end{equation}
两个方向取小值，是因为指令下发与状态回传均须可用。将附件对应接口参数代入，可得运输机—网关、运输机—中继和中继—网关的双向预算分别为122、116和126 dB。这些是设备参数推导的门限，不是某条实际航迹的通信余量。

设$D_{ab}(t)$为端点三维直线距离，单位km；频率$f$以MHz计，附件2.4 GHz对应2400 MHz。传播损耗为
\begin{equation}
L_{ab}^{path}(t)=32.45+20\log_{10}f+20\log_{10}D_{ab}(t)+L_{obs}\beta_{abt}.
\end{equation}
其中$L_{obs}=10$ dB，遮挡变量由三维视线与沿途DEM比较确定。系统损耗已在预算侧扣除，不能在传播损耗侧再次扣减。定义
\begin{equation}
A_{ab}(t)=\ind\{L_{ab}^{path}(t)\le L_{ab}^{max}\},\quad
M_{ab}^{link}(t)=L_{ab}^{max}-L_{ab}^{path}(t).
\end{equation}
遮挡增加损耗，但不必然导致链路中断；最终仍须比较预算门限。

\subsection{合法通信拓扑与全时段服务}
图\ref{fig:q3-topology}给出题目允许的两种通信方式。直连不可用时，可以经同一架中继完成接入和回传，不允许通过多架中继接力跳转。
\begin{figure}[htbp]\centering
\begin{tikzpicture}
\node[block,text width=2.5cm] (t) at(0,0){运输无人机};
\node[block,text width=2.5cm] (r) at(5,0){在服中继};
\node[block,text width=2.5cm] (g) at(10,0){固定网关G01};
\draw[edge,<->](t)--node[above,font=\small]{接入链路}(r);
\draw[edge,<->](r)--node[above,font=\small]{回传链路}(g);
\draw[edge,<->](t.south)--++(0,-1.1)--node[below,font=\small]{直连链路}(10,-1.6)--(g.south);
\node[font=\small,align=center] at(5,1.4){中继模式要求接入与回传同时可用\\拓扑示意，不表示实际点位或覆盖范围};
\end{tikzpicture}
\caption{直连与单中继通信拓扑示意}\label{fig:q3-topology}
\end{figure}

令$\mathcal T_r$为运输架次需要保障的全部飞行与交接时刻，$I_h^{srv}$为中继$h$的在服区间。对任意$t\in\mathcal T_r$，要求
\begin{equation}
A_{rG}(t)=1\quad\text{或}\quad
\exists h:\ t\in I_h^{srv},\ A_{rh}(t)=1,\ A_{hG}(t)=1.
\label{eq:continuous}
\end{equation}
输出时优先标记直连；非直连时选择一架满足条件的中继，保证每时刻只有一个实际保障指派。允许多个中继在物理上可用，但不能由此把同一运输区间同时计作多项必要保障任务。

\subsection{中继悬停、能量与资源占用}
候选悬停点$p$须在DEM内，海拔为$H_p^R=H_p^{DEM}+h_p^{AGL}$，且$0<h_p^{AGL}\le300$ m。中继出航巡航海拔还要覆盖终点悬停海拔：
\begin{equation}
H_{O01,p}^{c}=\max\left(\max_{c\in\mathcal C_{O01,p}}H_c+50,\ H_p^R\right).
\end{equation}
出航按爬升、水平巡航、必要下降计算，返航同理。中继从准备开始到建链完成的时刻依次为
\begin{align}
\tau_h^{dep}&=s_h+t_R^{prep},&
\tau_h^{arr}&=\tau_h^{dep}+t_{O01,p}^{fly},\\
\tau_h^{link}&=\tau_h^{arr}+t_R^{link},&
\tau_h^{ret}&=\tau_h^{end}+t_{p,O01}^{fly}.
\end{align}
只有在$[\tau_h^{link},\tau_h^{end}]$内才能提供保障。总能耗写为
\begin{align}
E_h^{fly}&=\frac{P_R^{cr}}{3600}\left(\frac{d_{O01,p}}{v_R^{cr}}+\frac{d_{p,O01}}{v_R^{cr}}\right)
+\frac{m_Rg_0(h_{O01,p}^{+}+h_{p,O01}^{+})}{3.6\times10^6\eta_R^{up}},\\
E_h&=E_h^{fly}+\frac{(P_R^{hover}+P_R^{com})(\tau_h^{end}-\tau_h^{arr})}{3600}
\le(1-\rho_R)E_R^{use}.\label{eq:relay-energy}
\end{align}
功率以kW计，故时间需除以3600。这里把建链期间按悬停及通信功率计入能源消耗，明确计费从抵达开始，而服务可用从建链完成开始。

中继实体机占用$[s_h,\tau_h^{ret}+t_R^{turn})$；能源组件占用任务及返航后充电区间。两架中继机和六组能源组件均按实体编号分配，在同一波次内也须逐一预占，不能让多个任务同时选中尚未更新状态的同一资源。

\subsection{候选覆盖与联合调度方法}
先将运输任务展开为分段线性三维轨迹，识别直连不足的时空位置，再从服务区邻域、航段附近和DEM高点形成候选悬停集合。依次剔除越界、回传不足和往返能量不可行的点。离散轨迹样本上的覆盖关系可用于快速筛选，但最终须按式\eqref{eq:continuous}验证整段时间。

联合候选方案采用硬截止优先组织运输批次，并根据中继最早建链时刻、最长可服务时长和能源再可用时刻调整波次。比较指标为一般物资迟到、联合完成时间、总能耗和总架次，其中
\begin{equation}
C_{max}^{joint}=\max\!\left(\max_r\tau_r^{ret},\max_h\tau_h^{ret}\right),\quad
E^{joint}=\sum_rE_r+\sum_hE_h.
\end{equation}
完成时间取最后返航，不包含此后的充电及周转。硬时限与连续通信必须先满足；不能用较少的中继能耗补偿已发生的通信中断。

\subsection{连续性验证与结果表达}
固定时间网格不能排除相邻样本之间的短时中断。验证时，应在每个线性运动区间给出传播损耗上界：距离项在区间端点求上界，地形项采用扫掠视线与经过像元的关系获得保守判定；无法判定的区间进一步细分。区间达到细分阈值仍未被证明可用时，应标记为未通过证明，不将其自动视为可行，也不直接断言实际已经中断。

后续需将认证区间与额外加密采样对照，并输出运输架次、阶段、起止时刻、直连或中继状态及实际中继架次。通信状态图、最小链路余量与联合资源甘特图分别回答“何时由谁保障”“距离门限有多少余量”“资源能否按时在位”。本阶段不报告尚未基于复核后轨迹得到的零中断、余量或中继数量结论。
